
Model-centric robustness concerns how the generator interprets, weighs, and integrates evidence once it appears in the prompt. Even with the same retrieved context, different generators can answer differently because they vary in instruction following, parametric prior strength, context utilisation, and calibration. Model-centric work therefore studies the internal and behavioral mechanisms that determine whether retrieved evidence is used faithfully or misleadingly.

\begin{table*}[t]
\centering
\small
\begin{tabularx}{\linewidth}{p{0.2\linewidth}p{0.2\linewidth}p{0.28\linewidth}X}
\toprule
\textbf{Problem} & \textbf{Representative work} & \textbf{Mechanism} & \textbf{Limitation or open issue} \\
\midrule
Parametric--retrieved conflict & \citetalias{tugOfWar2024} & Construct and analyze conflicts between internal and external knowledge & Conflict behavior depends on model, prompt, evidence coherence, and task framing. \\
Fact-level conflict modeling & \citetalias{faithfulRAG2025} & Extract parametric self-facts, align them with retrieved facts, and reason before answering & Requires reliable fact extraction and alignment; complex claims remain difficult. \\
Source-aware consolidation & \citetalias{astuteRAG2025} & Elicit internal knowledge and consolidate it with retrieved sources using reliability comparison & Training-free but still depends on LLM self-assessment and source-quality cues. \\
Hidden-state rectification & \citetalias{rrag008} & Locate layers where parametric suppression occurs and contrast contextual/non-contextual logits & Promising for model-centric control, but model internals and transferability remain open. \\
Context utilisation benchmarks & \citetalias{rrag032} & Test gold, conflicting, and irrelevant contexts across models & Benchmarks reveal whether a model uses context, but not always why a model chooses one source. \\
Real-document conflict & \citetalias{rrag112} & Study behavior under context-memory conflicts with real documents & Realistic evidence can be messier than synthetic conflicts and harder to label. \\
Source preference & \citetalias{rrag132} & Analyze which sources win under conflict & Source preference may encode bias, recency effects, or authority priors. \\
\bottomrule
\end{tabularx}
\caption{Representative model-centric work on knowledge conflict and contextual faithfulness.}
\label{tab:conflict-methods}
\end{table*}


\paragraph{Context Utilisation and Intermediate Reasoning}

One model-centric strategy is to force the generator to process evidence more explicitly before answering. Chain-of-Note asks the model to write notes for each retrieved passage, assess relevance, and abstain when evidence is insufficient \citep{chainOfNote2023}. This moves the generator from direct answer synthesis to a read--judge--answer pipeline. The value of this design is not simply better prompting; it externalises intermediate evidence judgments that can be inspected, scored, or used to trigger abstention.

Other approaches manipulate the context itself. Context compression and filtering reduce the amount of irrelevant information passed to the model, while context-utilisation benchmarks test whether models actually rely on gold evidence when it is present and ignore conflicting or irrelevant evidence when it is not \citep{rrag032}. The key distinction is between \emph{context availability} and \emph{context use}. A passage can be present in the prompt but ignored; a misleading passage can be unsupported but still dominate the answer. Thus, robust RAG must evaluate how models allocate attention, probability mass, and rationales across sources.

\paragraph{Parametric--Retrieved Knowledge Conflict}

RAG introduces a conflict structure that does not exist in standard closed-book generation. The model has parametric memory, the retriever supplies external evidence, and multiple retrieved passages may disagree. When these sources conflict, models can exhibit two opposite failures: over-compliance with false retrieved evidence or stubborn reliance on stale parametric knowledge. Tug-of-War constructs and studies internal--external conflict in retrieval-augmented language models \citep{tugOfWar2024}. Adaptive Chameleon and related work show that LLMs can be highly receptive to coherent counter-evidence but resistant under other conflict conditions \citep{rrag046}. Real-document conflict studies show that behavior under conflict is not only a synthetic benchmark artifact \citep{rrag112}.

Conflict-aware methods differ in granularity. FaithfulRAG extracts parametric self-facts, aligns them with retrieved facts, localises discrepancies, and uses a self-thinking step before generation \citep{faithfulRAG2025}. CoRect applies hidden-state and logit rectification to resolve conflicts without full model retraining \citep{rrag008}. Astute RAG elicits internal knowledge and consolidates it with retrieved passages through source-aware reliability comparison \citep{astuteRAG2025}. These methods share a principle: conflict should be represented explicitly before the final answer is produced.

\paragraph{From Faithfulness to Source Reliability}

Many faithful RAG methods implicitly assume that retrieved context should be trusted. Robust RAG cannot make this assumption. Retrieved context can be stale, poisoned, incomplete, or irrelevant; parametric knowledge can be outdated or hallucinated. The correct policy depends on source reliability. Source preference studies ask which sources win under conflict \citep{rrag132}; the \emph{Who's Who} evaluation of \citet{rrag110} studies behaviour under practical entity-level conflicts. These works suggest that future systems should represent not only claims but also provenance, timestamp, authority, and consistency. In a high-stakes setting, the most robust answer may not be a single synthesised fact; it may be a statement that exposes disagreement and explains why one source is more credible.

\paragraph{Citation, Attribution, and Verification}

RAG systems are often expected to cite sources, but citations do not guarantee faithfulness. A model can cite a passage that is topically related but does not support the generated claim, or cite the correct document while generating a claim from parametric memory. RAGTruth provides fine-grained hallucination annotations for RAG outputs \citep{ragTruth2024}. Related grounding methods train models to quote source spans or decompose external and parametric semantics, but these are not substitutes for claim-level citation verification \citep{rrag115,rrag206}. The key challenge is granularity. Document-level citation is too coarse when a passage contains multiple claims, qualifiers, or contradictions. Claim-level or sentence-level attribution is more informative but requires reliable entailment, alignment, and segmentation.

Verification can be pre-generation, post-generation, or iterative. Pre-generation verification asks whether the context is sufficient and trustworthy before answering. Post-generation verification checks whether each generated claim is supported by the evidence. Iterative verification alternates between generation and retrieval, using detected gaps or conflicts to form new queries. The limitation is cost: verification often requires additional model calls, retrieval calls, or entailment models. Robust RAG therefore needs cheap first-pass signals for routine cases and stronger verification for high-risk cases.

\paragraph{Limits of Model-Centric Fixes}

Model-centric methods can improve evidence use, but they cannot fully solve evidence unreliability. A generator cannot reliably infer provenance if the corpus does not preserve it. It cannot resolve a conflict if all sources are equally stale or poisoned. It cannot calibrate answerability if retrieval never exposes missing evidence. Conversely, evidence-side controls cannot guarantee faithful generation. This complementarity motivates the layered framework: evidence-centric controls shape the input distribution, model-centric controls shape context use, and system-centric controls decide when the remaining uncertainty is too high to answer.
